\chapter{天文与物理综合项目}
\chaptergoal{独立选择可行的天文或物理任务，建立原始数据与字段记录，制定数据划分和基线方案，选择两种具有明确理由的模型，完成可信的验证与一次最终评价，并提交能够复现图表和结论的项目报告。}

前十六章提供了模型与工具。本章将它们放入实际研究工作的顺序中：不是先想好用哪一个高性能算法，再寻找数据让它得到高分，而是先明确问题和可获得的信息，再决定基线、候选模型、评价方式和结论范围。课程内的八学时用于选题、方案检查、报告讨论和答辩，完整的数据处理与分析必须安排一定课外时间。
\section{选题、任务定义与可行性}

\textbf{综合项目}是一个有明确起点与交付结果的小型研究任务，不要求发明新算法。前面已介绍的数据对象在这里按任务重新汇总；不熟悉某一数据时，应先返回首次介绍，而不是从下载函数开始猜它的含义。

\begin{table}[htbp]\centering\small
\begin{tabularx}{\linewidth}{p{0.25\linewidth}Xp{0.17\linewidth}}
\toprule 数据与对象 & 本项目主要输入与目标 & 首次介绍 \\\midrule
造父变星目录（OGLE） & 周期预测平均I波段视星等；每行一颗恒星 & 第\ref{sec:ogle-intro}节\\
天体观测目录（SDSS） & 测光特征预测参考光谱类别；按天体分组 & 第\ref{sec:sdss-intro}节\\
天体测量目录（Gaia） & 位置、视差和自行寻找候选群体；无直接训练类别 & 第\ref{sec:gaia-intro}节\\
翼型自噪声（Airfoil） & 实验条件预测声压级；每行一组测量条件 & 第\ref{sec:airfoil-intro}节\\
声呐回波（Sonar） & 60个回波特征区分两类实验目标 & 第\ref{sec:sonar-intro}节\\
超导材料记录 & 组成描述特征预测临界温度；注意材料重复 & 第\ref{sec:superconduct-intro}节\\
\bottomrule
\end{tabularx}
\caption{综合项目数据入口。表中是任务设计，不是已经全部完成真实数据复算的结果目录。}
\end{table}

把“我想使用随机森林”改写为可执行问题，可以是：“使用翼型实验中的五项条件，预测尚未用于训练的测量记录的声压级，并比较线性模型与随机森林在相同验证记录上的MAE。”这句话规定了输入、目标、对象和比较方式；如果进一步要预测全新工况，就必须另行定义工况划分，不能沿用随机留出而不作解释。

一个好的本科项目应有清楚的研究问题、可取得的数据、可检查的评价标准和有限的工作量。建议从以下任务中选择：OGLE 造父变星周期--星等回归；SDSS 测光类别预测；Gaia 星团邻域候选识别；Airfoil 风洞声压级预测；Sonar 声呐回波分类。具备较强专业与编程基础的学生，还可尝试超导材料临界温度预测。

不必在一次项目中同时使用所有算法。推荐从一个简单基线和两种有科学或方法依据的候选模型开始，保持比较清晰。任务定义至少包含样本单位、可用输入、目标及其参考来源、应用对象、评价指标和主要风险。教师审核时应特别确认数据能否获得，标签是否真实有效，是否存在重复对象或同工况相关记录。

\fig{ch17_project_flow.pdf}{科学机器学习项目的七个核心阶段：问题、数据、划分、基线、模型比较、独立评价及科学解释。}

\subsection{本章提供两条完整、性质不同的实践路线}
为了避免综合项目只停留在选题清单，本章加入两项已经执行的示范：一项使用真实OGLE大麦哲伦云基频造父变星目录的\textbf{固定90条节选}，研究周期与平均 $I$波段视星等关系；另一项使用\textbf{公开生成规则的物理标定模拟数据}，研究在未见实验批次上预测传感器电压。前者检验真实字段和经验关系，后者可以控制批次偏移、噪声和真实函数，从而检查方法是否恢复已知结构。两者不能互相冒充。

真实目录节选不是完整巡天样本，也不是随机代表全体大麦哲伦云造父变星的样本。选取规则事先固定为源文件前90条记录，没有依据预测分数挑选对象。它足够用于一项从读表到评价的课堂实战，却不足以发表新的周期—光度标定或推断总体完备性。完整目录项目可以在掌握流程后扩展，届时需重新审查质量、选择效应和评价。

物理模拟则不是某台真实仪器的性能测试。其价值在于知道哪些项真实存在，可以检查模型缺少二次项或忽略批次结构时会出现什么错误。把模拟作为一个透明的科学计算实验，比把无法获得的观测数据换成未标注的生成数字更可靠。

\subsection{用任务合同限定比较的范围}
每个项目开始时写一段简短的“任务合同”：样本单位、部署时可获得输入、预测目标、目标单位、预期新样本、主要指标、数据范围和资源预算。例如标定项目的合同是“对训练时未出现的实验批次，利用输入激励和温度预测输出电压，以按批次留出的RMSE评价”。它排除了把批次编号作为连续物理预测量，也明确随机拆行不是最终评价目标。

模型选择只比较一个简单基线与两三种有明确理由的方法。先利用任务合同判断一个新建议是否回答原问题，再决定是否增加模型。不断试更多算法直到某个测试分数变高，不是独立项目的充分分析。

\section{数据来源、字段与质量检查}

\textbf{数据版本}指本项目实际使用的那一批文件和其整理状态。相同下载地址可能随项目更新而返回不同内容，因此保存获取日期、原始文件摘要和处理脚本，才能准确说明“同一数据”。\textbf{文件摘要}是从文件字节计算的校验值，用于检查副本是否一致，不用于证明数据科学上正确。

先抽查若干行：读出一条记录的对象、输入和目标，并把单位写成完整句子。再统计缺失和重复。只有在这一步能够解释清楚每一列，才能继续模型训练；“文件能被Pandas读入”不是字段含义已经核对的证据。

优先使用官方数据源或可追溯的公开整理数据。每个项目创建 \texttt{data\_dictionary.csv}，记录列名、类型、单位、含义、缺失编码、角色和来源。如果数据文件不能合法再分发，需提供下载链接、检索条件和整理脚本。对含对象标识的表格，要在任何连接与分组前确认标识字段没有浮点精度损失。

对天文目录，应记录发布版本、查询范围、质量门槛、交叉匹配原则、多个观测对应的对象身份。对物理实验，应记录测量条件、仪器或工况编码、重复实验关系，以及测量误差或质量标记。任何被删去的记录都应能够解释原因，不能把所有不利于模型的样本统称为“异常值”。

\begin{table}[htbp]\centering\small
\caption{专业项目数据字典的最低要求}\begin{tabularx}{\linewidth}{p{0.19\linewidth}X}\toprule
字段&需要回答的具体问题\\\midrule
数据来源&机构/项目、版本、获取日期、链接与许可\\
样本单位&一个对象、一次观测还是一次实验工况？\\
字段角色&输入、目标、身份、分组、质量标记还是外部检查字段？\\
单位/缺失&单位与数值范围；缺失和未检出的区别\\
筛选条件&删除或保留记录的可复现规则\\
关联关系&是否重复对象、相关工况或同材料配方？\\\bottomrule
\end{tabularx}\end{table}

\subsection{真实天文示范：从目录字段开始}
OGLE是长期开展光度监测的天文项目，本例来源为其公开变星目录中的大麦哲伦云经典造父变星基频文件 \texttt{cepF.dat}。造父变星亮度具有周期变化；基频与其他脉动模式应分别处理。文件说明给出源编号、强度平均 $I,V$波段视星等、周期等字段。本例提取前90条的编号、$I$、$V$和周期，并保留原始零起始行号。[13]

\begin{center}\small
\begin{tabularx}{\linewidth}{p{0.24\linewidth}p{0.15\linewidth}X}\toprule
字段&单位&含义与用途\\\midrule
\texttt{source\_id}&无&目录对象标识，按字符串保存，不进入回归输入。\\
\texttt{source\_row}&无&原文件行位置，用于核对节选，不用于预测。\\
\texttt{period\_days}&天&目录给出的脉动周期，本例的预测输入。\\
\texttt{I\_mag}&mag&强度平均 $I$波段视星等，本例目标。\\
\texttt{V\_mag}&mag&保留供字段核对和扩展检查，本次不作为输入。\\\bottomrule
\end{tabularx}\end{center}

空白单元表示源文件缺测，不替换成0。本例90条中3条缺少 $I$目标，不能用于当前监督评价，保留其排除记录；其他87条周期均为正，可计算对数。某些记录缺少 $V$，但本例只用周期预测 $I$，没有必要因为一个未使用字段缺失而删除它们。

87条可用记录的周期约1.263至28.374天。第0条编号为OGLE-LMC-CEP-0002，周期3.118149天，$I=15.691$mag；它既不是人造样本，也不是某个“理想关系”上的拟合点。节选和源字段说明一同保存，附录提供联网后对源文件前90条重新核验的脚本。

\subsection{物理示范：明确给出生成机制}
模拟标定有90个实验批次，每批8条记录，共720条。每批温度 $T$ 从15至35摄氏度范围抽取并在批内固定，输入激励 $x$ 在0至10范围生成。无噪声响应为
\begin{equation}
f(x,T)=0.5+1.8x+0.06x^2+0.02(T-25)x,
\label{eq:project-calibration}
\end{equation}
输出单位设为伏特，输入单位设为归一化激励单位。每个批次还叠加一个标准差0.12伏的共同零点偏移，每条记录再叠加标准差0.08伏的独立高斯测量噪声。生成种子为170，全部原始生成代码包含在工程中。

分析表包含记录编号、批次编号、输入激励、温度和观测输出。批次编号用于分组而不是数值特征；共同偏移不提供给模型，因为真实应用中它通常未知。知道生成机制是教师验证分析的条件，但学生建模时只使用声明可获得的两个输入，不能把无噪声真值或隐藏偏移混入预测特征。

数据检查包括唯一记录数、每批样本数、温度在批内是否一致、输入和输出是否有限、字段单位是否一致。人工生成数据也需要检查：程序写错单位或把目标列复制进输入，同样会破坏教学实验。

\section{划分依据与基线建立}
项目必须先根据目标应用设计数据划分，再开展复杂模型比较。SDSS 的多次观测可能要求按唯一对象分组；Airfoil 若存在可以识别的重复工况，应讨论是否需要按工况留出；材料任务若要预测全新材料族，就需要相应分组评价，而不是只随机拆分同族记录。没有可靠分组字段时，应明确局限，不能凭空制造分组身份。

建立简单基线是必做工作。分类可采用多数类预测或逻辑回归；回归可采用训练集均值或线性回归；聚类和候选任务则先画出原始特征与简单几何选择结果。基线的目的不是获得理想最终结果，而是给复杂方法提供可以解释的参照。

\subsection{天文示范的对象划分}
87条可用造父变星记录中每个源编号唯一。固定种子42，将69条作为开发数据、18条最终保留。开发内部使用5折、打乱种子7；这个划分用于课堂预测比较，不能代表完整巡天的空间外推。由于节选本身来自固定文件前缀，讨论总体代表性时必须保留这一限制。

常数基线从每个训练折的 $I$目标计算均值，不使用验证目标。在星等空间直接以周期作线性回归，是一个可解释但可能过于简单的候选；使用 $\log_{10}P$作输入则检验常用经验表示。对数周期是任务提出前指定的特征变换，不是看测试曲线后临时试出的修正。

这里预测的是\textbf{同一选定区域与模式样本中的平均视星等}，没有从数据中获得每个对象的独立距离、消光或金属丰度校正。因此回归截距不能直接当作绝对星等零点，拟合残差也不能全部归为测量误差。先明确目标，才有资格解释系数。

\subsection{物理示范为什么按批次划分}
同一批次8条记录共享温度与零点偏移。若随机拆行，训练通常已经见过该批的其他记录，验证更接近“已见批次的新测量”。本项目希望处理新批次，所以先按90个批次留出18批作为最终测试，共144条；其余72批、576条用于开发。开发内部使用5折 \texttt{GroupKFold}，同批记录不跨训练与验证。

\begin{lstlisting}[caption={标定项目先保留整个批次，再做分组交叉验证}]
from sklearn.model_selection import GroupShuffleSplit, GroupKFold
X = data[['input_level','temperature_C']]
y = data['response_V']
groups = data['run_id']
split = GroupShuffleSplit(n_splits=1, test_size=.2, random_state=42)
dev, test = next(split.split(X, y, groups))
cv = GroupKFold(n_splits=5)
# Pass groups.iloc[dev] to the development cross-validation.
\end{lstlisting}

整体均值基线不利用输入；线性模型加入输入和温度；二次特征模型允许曲率与交互；森林提供一个不同结构的非线性参照。四个模型可以回答逐步具体的问题：输入是否有信息，线性是否足够，已知形式的低阶非线性是否改善，以及更灵活的树模型是否必要。

\section{模型比较与有限调参}
模型比较应尽量控制变量：固定输入定义、开发范围和评价目标，只改变有明确理由的模型与有限超参数。对回归可选线性回归、随机森林和梯度提升；对分类可选逻辑回归、森林和 MLP；对无监督任务可选 K-means 与 DBSCAN，并重点比较结构假设和结果稳定性。不能为了分数最高而进行无限制的组合搜索。

使用 \texttt{Pipeline} 管理标准化、填补与模型，使交叉验证中每折只利用该折训练记录拟合预处理。对于需要类别不平衡处理的任务，任何重采样或权重策略也必须在训练范围内完成。报告中应列出候选模型及其主要参数，并说明为什么尝试这些方法。

\begin{lstlisting}[caption={综合项目中保存统一模型比较记录的示例}]
import pandas as pd
from sklearn.model_selection import cross_validate

rows = []
for name, estimator in candidate_models.items():
    result = cross_validate(estimator, X_dev, y_dev,
                            cv=cv, scoring=scoring, n_jobs=1)
    rows.append({"model": name,
                 "mean_score": result["test_score"].mean(),
                 "std_score": result["test_score"].std()})
pd.DataFrame(rows).to_csv("comparison.csv", index=False)
\end{lstlisting}

示例假设 \texttt{candidate\_models}、\texttt{cv} 和 \texttt{scoring} 已在项目脚本中定义。这是一段承接项目定义的代码片段，不是单独运行的完整程序。工程中的\texttt{scripts/run\_science\_csv.py}提供可执行的表格任务入口，可先用\texttt{--help}检查所需参数；使用专业数据前仍要自行核对目标列和分组字段。不要在尚未获取真实目录时，给模板插入模拟数据后宣称完成了专业验证。

\subsection{天文示范的三个预先规定模型}
在69条开发记录的相同五折上，比较均值、原始周期线性模型和对数周期线性模型，得到以下RMSE，单位均为mag。每个模型在各训练折重新拟合；表中的折间标准差使用分母为折数减一的样本标准差，见对应脚本。对数变换没有需要从全体样本学习的参数。
\begin{center}\small
\begin{tabular}{lrr}\toprule
模型&开发五折RMSE均值&折间标准差\\\midrule
训练均值&0.7524&0.1600\\
$I=a+bP$&0.4201&0.1110\\
$I=a+b\log_{10}P$&0.2217&0.0882\\\bottomrule
\end{tabular}\end{center}

本次对数周期模型在开发均值上最好，因此按预定主要指标选它。标准差较大说明这个小节选对具体划分敏感，不能把0.2217写成完整目录的稳定精度。原始周期模型已比常数好，表明周期有预测信息；对数表示进一步改善，说明特征表达可以比盲目增加算法复杂度更重要。

在全部69条开发记录上重拟合后，得到
\[
\widehat I=16.9021-2.8806\log_{10}P.
\]
这两个系数属于当前节选、当前划分和未作额外消光等校正的经验关系。它们用于重现课堂预测，不作为新的天文标定结果发布。

\subsection{物理示范的开发比较}
二次模型采用 \texttt{PolynomialFeatures(degree=2, include\_bias=False)}，随后标准化并用 \texttt{Ridge(alpha=0.01)}拟合；森林固定200棵、最小叶样本数3。所有步骤在每个开发训练折内重建，未使用最终18批的输出选择配置。
\begin{center}\small
\begin{tabular}{lrr}\toprule
模型&分组五折RMSE均值（伏）&折间标准差（伏）\\\midrule
训练均值&6.7594&0.4954\\
线性模型&0.5869&0.0619\\
二次特征加Ridge&0.1440&0.0170\\
随机森林&0.3307&0.0621\\\bottomrule
\end{tabular}\end{center}

这个模拟的真关系包含二次项和温度交互，所以低阶特征模型有适合的表示能力。它优于当前森林，不是因为所有物理任务都应采用二次回归，而是因为本实验结构恰好支持这种模型。森林无需预先写出关系，却需要用许多分段近似平滑响应；在有限样本下，这种灵活性不一定带来更低误差。

\subsection{比较结果后应如何停止选择}
项目可以在预定候选比较后选定二次模型，不必为了压低到更多小数位继续使用测试集试模型。若开发残差仍有明确结构，可以事先再设置一轮有限研究，但须记录新假设与所有候选。方案包含输入、预处理、模型、超参数和评价规则，不能只冻结算法名称却继续改特征。

本章示范选择后，在全部开发数据上重拟合一次，然后才打开最终保留结果。最终测试不用于重新决定周期形式、二次阶数或森林叶大小。这样，读者能明确区分探索结果与最终检查。

\section{最终评价与错误分析}
监督预测项目在开发阶段完成模型和超参数选择，再使用预先保留的独立测试部分做一次最终评价。评价应同时给出总体指标、样本数量及至少一种错误/分组诊断。若模型在某类对象上存在明显偏差，不应只保留汇总高分。

对于 Gaia 等无监督候选项目，不能生搬硬套分类准确率。需要使用聚类稳定性、参考样本、外部颜色--星等关系、天体测量质量和人工核查等证据，并明确这些证据在选择方案时是否已被使用。若有可靠的外部成员标签，可作事后对照，但其来源和不确定性也应写清楚。

\subsection{真实目录节选的最终保留结果}
对18条最终保留的造父变星记录，对数周期模型得到RMSE约0.1492mag、MAE约0.1248mag，平均预测减目标约0.0555mag。这些值低于开发折均值，不必强行解释成模型突然变好；两部分包含不同对象，小样本分配与极端残差足以造成差异。

\fig{v3_project_ogle_fit.pdf}{真实OGLE目录固定90条节选中的87条可用记录。曲线仅由69条开发记录拟合，18条最终保留记录另用符号标记；星等轴反向显示，数值较小表示更亮。}
\fig{v3_project_ogle_residual.pdf}{18条最终保留天体的预测减目录 $I$星等残差。该图是选定方案后的检查，不用于再筛除误差大的对象或重新选择模型。}

最大绝对残差约0.2935mag，对象为OGLE-LMC-CEP-0048，周期约5.5482天，目录 $I=14.465$，预测约14.7585。它比模型预测更亮，但仅凭这个残差不能判断原因，更不能把该点从评价中删除以改善RMSE。需要回到原目录、光变曲线、拥挤与消光等信息核验。本教材节选没有这些全部信息，因此将其作为待检查的误差对象。

\subsection{新实验批次上的最终结果}
二次Ridge在576条开发记录上重拟合后，对18个未见批次的144条测试记录，RMSE为0.1473伏，MAE为0.1170伏。这个水平接近模拟中未提供给模型的批次偏移与测量噪声共同造成的尺度，但不应通过一个数就声称达到精确理论下限；估计误差、有限样本和残差相关性仍有贡献。

\fig{v3_ch17_predictions.pdf}{人工物理标定项目的最终批次保留预测。对象为模拟实验，不能据此宣称某台真实传感器具有相同精度。}

按批次汇总时，批次44的平均残差约0.2968伏，批次86约 $-0.2517$伏；其他批次较小。这与同批共同零点偏移的生成机制相符：模型没有获得隐藏偏移，可能在一整批上同向偏离。若按行看所有144条残差而忽略批次，就会低估这类系统变化的重要性。

\subsection{错误分析不是测试后的重新调参}
可以用最终残差说明模型的限制，提出未来需要新增温漂校正或批次参照测量的研究方向。若真的据此改变模型，就需要新的独立评价或明确称为后续探索，不能再次把同18批当作未参与设计的最终测试。论文与课程报告都应记录这个时间顺序。

对于本节两项示范，完整预测表保存对象或批次标识，能够从任意一个指标追溯到对应记录。没有仅保留最漂亮的散点图，也没有用删除大误差点后的分数替代预定评价。

\section{科学解释与结论边界}
科学解释的第一步是检查结果与已知物理规律是否一致，第二步是分析无法解释或预测不好的区域。对于 OGLE，周期--星等关系是否随样本筛选改变？对于 SDSS，分类错误是否集中在暗弱天体或颜色重叠区域？对于 Airfoil，高频或高攻角工况是否具有较大残差？对于 Gaia，聚类候选是否在天空位置及颜色--星等图中具有合理一致性？

\textbf{重要的是不做超过数据证据的结论。} 对训练范围内的相关关系、对新观测对象的预测能力和对物理机制的论证，分别有不同的证据标准。专业项目可以清楚地给出一个局限明确但可信的结果，而不必宣称发现新的宇宙规律。

\subsection{天文示范能够支持与不能支持的陈述}
本节可以支持：“在公开目录前90条的可用基频样本中，对数周期模型比预先列出的均值和原周期线性候选具有更低开发RMSE；按固定划分保留的18条记录得到约0.149mag的预测RMSE。”这包括来源、样本、候选、评价和数值。

它不能直接支持：“我们获得了大麦哲伦云全体造父变星的最新周期—光度定律。”原因包括节选不代表完整总体，目标是视星等而非独立绝对光度，未系统处理消光、几何深度、金属丰度、拥挤、模式混入及选择效应。必要限制不是为结果辩护，而是定义研究究竟做了什么。

已有经验关系为采用对数周期提供了背景，本课程的工作是练习从真实数据建立、比较与评价模型，而不是声称首次发现这一关系。扩展到科学研究时，应检索相应文献、完整处理样本和误差，并与既有标定比较。

\subsection{物理模拟能够验证什么}
模拟中知道真实函数与噪声分量，因此可以检查算法能否表达已知曲率、分组是否避免共享偏移泄漏，以及残差是否显示预期的批次相关。这是受控计算实验。它检验方法在明确条件下的行为，不提供现实仪器所有误差源的真实性保证。

若要转向真实实验，首先核对传感器饱和、迟滞、频率响应、温度测量误差和标定参照等因素。缺少的效应不能因为模拟没有加入就视为不存在。真实数据也可能不存在相同的二次形式，因此应重新比较基线而不是直接固定本章获胜模型。

\subsection{把下一步研究写成可检验问题}
一个有用的后续问题应指出需要什么新数据或设计。例如天文方向可比较不同天区、模式或消光处理下的残差，物理方向可增加独立批次校准点，检验零点修正是否改善新批次误差。不要只写“采用更先进算法提高精度”；新方法应对应已经观察到的具体不足。

若用SDSS、Gaia、Sonar或Airfoil开展自己的项目，仍沿用这里的证据层级：数据含义、对象划分、合理基线、有限比较、真实输出、误差来源、结论范围。模型名称可以变化，研究逻辑不应消失。

\section{代码、报告与答辩}

\textbf{可复现}在本项目中的最低含义是：另一位同学得到相同数据、程序和环境说明后，能按记录的步骤重新生成核心表图，并理解数值差异可能来自哪里。它不等于只上传一个模型文件，也不等于结果在任何机器和所有软件版本中都逐字节相同。

可以先写一份简短的结果叙述模板：“在规定的开发样本中，我们比较了……；使用……划分与指标选择配置；冻结配置后，在……条保留记录上得到……；主要错误集中于……；尚不能据此判断……。”空缺处必须由实际运行和证据填写，不能从教科书示例的分数直接拷贝。

最终提交包含：报告、一个从干净环境可顺序运行的 Notebook、必要数据或获取脚本、字段字典、环境信息、划分方案、模型配置、原始比较表及生成图表的代码。报告建议依次安排背景与问题、数据、方法、结果、讨论与结论、可复现信息。

\begin{table}[htbp]\centering\small
\caption{综合项目评分建议，作为本章项目内部评价权重}\begin{tabularx}{\linewidth}{Xr}\toprule
评价维度&分值比例\\\midrule
问题定义与数据来源完整性&20\%\\
基线、模型选择与比较方法&20\%\\
数据划分与评价设计&25\%\\
错误分析、科学解释与结论边界&25\%\\
代码可复现性与报告表达&10\%\\\bottomrule
\end{tabularx}\end{table}

答辩至少回答三个问题：为什么把专业问题设计成当前的任务？为什么当前评价能够代表所关心的新样本？如果更换观测条件、仪器或材料族，模型还有哪些需要重新检验的地方？学生应准备一幅最重要的结果图和一幅最重要的错误或局限图，能够不依赖 PPT 朗读解释其意义。

\subsection{一份能够被另一位同学复现的提交}
工程根目录保留说明文件，写清运行顺序与必要依赖；数据目录含原始或经过许可的节选、字段字典和来源；划分文件使用稳定对象标识；结果目录保存预测表、指标和图表；脚本或Notebook能从已说明输入重新运行。随机数据还必须保存生成规则与种子，真实目录不能用模拟文件同名覆盖。

报告正文不需要逐行粘贴代码，而应解释关键选择。代码负责精确执行，正文负责说明为什么这样执行。一个读者应能从报告中的某张表找到生成脚本，也能从脚本输出回到相应段落。图表一旦重新生成，检查单位、记录数和模型配置是否与文字一致。

在工程根目录，物理模拟与天文节选项目分别运行：
\begin{lstlisting}[language=bash]
python experiments/ch17.py
python experiments/project_ogle.py
\end{lstlisting}
它们都从已随书保存的数据或明确生成过程开始，不依赖隐藏的个人路径。外部完整数据扩展另有获取脚本，失败时清楚报错，不用替代数据悄悄继续。

\subsection{答辩应能回答的关键问题}
说明一条样本是什么、目标从哪里来、部署时是否能得到全部输入；解释为什么选择当前划分而不是随机拆行；指出基线的含义；给出一项模型比较的实际证据；挑一个失败对象解释目前知道和不知道什么；说明结论在哪些范围外不能直接使用。能够回答这些问题，比记住一个最佳参数组合更能体现独立建模能力。

对于未完成的项目环节，应在提交中准确标明。数据下载脚本存在不等于已运行数据分析，算法执行成功不等于科学假设得到验证。课程可以评价一个诚实、分析充分但结果没有改善的项目，不应鼓励为了高分而删掉不利实验。

\subsection{本书方法如何组成一个完整工作过程}
第1章帮助明确任务，第2章建立可信数据表，第3章理解训练目标与优化，第4章规定评价。第5至12章提供可比较的预测方法，第13至15章处理表示、结构和异常，第16章把数值结果转为有边界的解释。本章将它们连接起来，形成从问题到证据的完整过程。

完成课程后，面对一份新数据，最先问的应是“我要预测或发现什么，哪些信息可以使用，怎样判断结果有效”，而不是“哪个最新模型可以直接调用”。模型能力来自数学表达与数据支持，科研可信度则来自整个分析过程。

\chapterpractice{从五个标准专业项目中选一项；提交一页选题说明、字段字典、数据划分与基线；在开发数据上比较两种候选模型；冻结方案后完成最终评价；提交报告、可运行代码和答辩材料。提高项目可使用超导材料数据，但必须处理材料分组与外推定义。}
\chapterexercise
\begin{enumerate}
\item 把“利用测光预测天体类别”写成可执行的研究任务定义，列出输入、标签来源、对象分组及评价目标。
\item 如果一份实验表格缺少工况身份，怎样诚实描述随机划分的适用范围？
\item 为什么期末项目不应仅按单次最高准确率评分？
\item 对同一批数据做多次参数搜索之后，如何保护最终测试数据的独立性？
\item 设计一张综合项目复现清单，使另一位同学能够重新生成核心图表。
\end{enumerate}

\section*{深化练习与方法检验}
本组题强调解释、计算与验证；D编号在本章内独立计数。计算题应写出所用定义和中间步骤，实验题保留实际输出，包括不支持预期的结果。
\begin{enumerate}[label=D\arabic*.]
\item 复算OGLE节选项目，逐步解释90条原始记录、87条可用记录、69条开发记录和18条保留记录之间的关系。V波段缺失是否应在当前模型中自动导致删行？
\item 对数周期模型的保留RMSE约0.149mag，能否据此报告对大麦哲伦云全部造父变星的精度？写一句包含来源、范围、评价和限制的结果陈述。
\item 物理模拟有90个实验批次，每批8条记录。为什么随机拆720行与按批次留出18组测试不是同一问题？已知真函数应怎样用于教学而不泄漏实际建模？
\item 从该工程选择一张结果表，写出完整追溯链：正文结论、结果文件、生成脚本、输入数据、划分、模型选择和最终评价。另一位同学在哪一步最容易误把开发结果当作测试结果？
\end{enumerate}
